\chapter{固定翼飞行器着陆控制系统组成与需求分析}

\section{固定翼飞行器着陆任务场景与系统组成}

固定翼飞行器着陆阶段在其整个飞行过程当中，属于风险最为集中且操纵要求颇高的阶段范畴。依据本文仿真平台的运行逻辑剖析可知，飞行器的着陆过程能够被划分成进近、拉平、接地区域、滑跑这四个阶段。进近阶段的主要任务是实现下滑轨迹的维持、速度的有效控制、跑道的精准对准；拉平阶段主要是降低下沉速率并且对俯仰姿态加以调整，让飞行器平稳地向地面靠近；接地区域着重关注接地高度、姿态和侧偏误差；滑跑阶段重点对接地后的方向保持与状态收敛展开分析。着陆后段飞行高度较低、修正时间较短且对环境扰动较为敏感，因此此阶段既要求控制系统具备稳定可靠性，又要求操纵者能够及时把握关键状态的变化情况并做出正确的决策。

就系统组成而言，本研究的固定翼飞行器着陆控制系统主要涵盖感知层、控制层、人机交互层和安全保护层这四个部分。感知层的作用是收集离地高度、速度、下沉率、俯仰角、滚转角、航向角、侧偏量以及地形高程等状态信息，它是实现闭环控制的基础。控制层会按照飞行状态和目标约束，实时产生油门、升降舵、副翼、方向舵和襟翼指令，以实现纵向和侧向的修正。人机交互层负责把关键状态、风险提示和建议操纵信息反馈给操纵者。安全保护层会依据阈值判断是否需要限幅、约束或者复飞。由此可知，着陆控制系统并非单纯的自动控制模块，而是一个将状态获取、控制决策、信息提示和保护机制融合在一起的综合系统。

\section{着陆控制关键参数与影响因素}

在固定翼飞行器着陆控制中，关键参数包括目标速度、最小速度与最大速度、下滑角、拉平高度、拉平距离、接地高度、最大下沉率、最大滚转角、最大俯角、最小俯角、最大侧偏量以及跑道捕获与提交边界等。其中，速度对着陆后段飞行器是否有足够升力和操纵裕度起决定作用；下滑角能体现整体下降的趋势；拉平高度和拉平距离决定着阶段切换的时机；下沉率、滚转角、俯仰角和侧偏量会直接对接地品质和安全性产生影响。速度过低会导致飞行器操纵迟钝，下沉率过大会明显增强接地冲击，滚转或侧偏过大可能导致偏离跑道中心线或接地姿态不正常。

影响着陆控制效果的因素主要可分为三类。第一类是环境因素，包括地形起伏、局部高程变化和近地扰动等。当外部环境发生变化时，飞行器相对地面的离地高度变化规律也会随之直接改变。第二类是飞行器本体因素，主要体现为机体姿态响应能力和起落架缓冲性能。着陆控制既要关注飞行轨迹，同时也要考虑接地载荷是否处在安全范围之内。第三类是操纵因素，即操纵输入是否平稳、修正是否及时以及末段判断是否准确。在复杂地形或近地阶段，缺乏有效提示会导致操纵者出现修正滞后或过度修正的情况，进而影响接地质量。

\section{研究对象参数整理与取值依据}

本文把小型固定翼飞行器着陆控制系统作为研究对象，运用典型参数建模的方法展开分析，并非针对某一具体实机做完全复现。按照当前平台的设置，研究对象的初始离地高度为120 ft，初始速度为72 kts，初始下降率为-250 fpm，跑道距离为2200 ft，目标接地点在跑道阈值后120 ft处，跑道捕获距离为320 ft，跑道提交距离为90 ft，跑道提交高度为22 ft。对于近地约束，还设定了拉平阶段最大侧偏28 ft、最大航向误差6°、最大滚转10°、接地提交阶段最大侧偏45 ft、最大航向误差10°、最大滚转12°。

这些参数的取值依照"典型、合理、可仿真"的原则。速度参数按照小型固定翼飞行器末段进近的通常范围确定，既要防止失速风险，又要抑制接地冲击。高度和距离参数主要用于区分进近、拉平和接地区域，以便系统依照不同阶段调用不同控制策略。横向约束参数用于保证飞行器在接近目标接地点时具备更严格的跑道对准能力。

\section{人机交互强化需求分析}

传统着陆控制研究通常更多地聚焦于飞行器能否沿着目标轨迹稳定下降，然而对于操纵者在着陆过程中的感知、判断和干预能力却没有给予足够的考量。事实上，着陆阶段属于典型的人机协同过程。结合现有平台界面能够看到，系统已经可以实时显示时间、距离、离地高度、速度、下沉率、俯仰角、滚转角、航向角、侧偏量和控制指令，并且能够输出"建议""告警""触发原因"等信息。这说明系统已然具备了进行人机交互强化研究的基础。

人机交互强化需求主要呈现于三个方面。第一是状态提示需求，即系统需要及时且清晰地向操纵者展示关键状态量及其变化趋势，帮助操纵者判定当前所处的飞行阶段和着陆风险。第二是风险告警需求，当系统检测出速度过低、下沉率过大、滚转超限、侧偏过大或者跑道对准不足时，应当及时发出告警，让操纵者能够提前进行修正。第三是操纵辅助与约束需求，即系统在高风险阶段不仅要进行提示，必要时还要通过限幅、约束或者复飞触发来防止误操作扩大风险。固定翼飞行器着陆控制系统的需求不仅仅是实现基本着陆，而是要在满足轨迹控制和接地安全的条件下，借助人机交互来增强操纵稳定性、风险感知能力和系统整体安全性。
